\section*{Referências}\addcontentsline{toc}{section}{Referências}

As referências devem ser apresentadas em ordem alfabética. Só podem ser
inseridas nas referências os documentos citados ao longo do artigo.
Todos os documentos citados obrigatoriamente têm que estar inseridos nas
referências. A seguir são apresentados alguns exemplos de referências
bibliográficas. Destaca-se que deve ser seguida a norma da ABNT.

\noindent\textbf{[parte de um documento:]}

AMADO, Gilles. Coesão organizacional e ilusão coletiva. In: MOTTA,
Fernando C. P.; FREITAS, Maria E. (Org.). \textbf{Vida psíquica e
organização}. Rio de Janeiro: FGV, 2000. p. 103-115.

\noindent\textbf{[trabalho acadêmico ou monografia (TCC/Estágio, especialização, dissertação, tese):]}

AMBONI, Narcisa F. \textbf{Estratégias organizacionais}: um estudo de
multicasos em sistemas universitários federais das capitais da região
sul do país. 1995. 143 f. Dissertação (Mestrado em Administração) -
Curso de Pós-Graduação em Administração, Universidade Federal de Santa
Catarina, Florianópolis.

\noindent\textbf{[norma técnica:]}

ASSOCIAÇÃO BRASILEIRA DE NORMAS TÉCNICAS. \textbf{NBR 6023}: informação
e documentação: referências - elaboração. Rio de Janeiro, 2002a. 24 p.

ASSOCIAÇÃO BRASILEIRA DE NORMAS TÉCNICAS. \textbf{NBR 10520}: informação
e documentação: citações em documentos: apresentação. Rio de Janeiro,
2002b. 7 p.

\noindent\textbf{[livro:]}

BASTOS, Lília R.; PAIXÃO, Lyra; FERNANDES, Lúcia M. \textbf{Manual para
a elaboração de projetos e relatórios de pesquisa, teses e
dissertações}. Rio de Janeiro: Zahar, 1979.

\noindent\textbf{[trabalho acadêmico ou monografia (TCC/Estágio, especialização, dissertação, tese):]}

BRUXEL, Jorge L. \textbf{Definição de um interpretador para a linguagem
Portugol, utilizando gramática de atributos}. 1996. 77 f. Trabalho de
Conclusão de Curso (Bacharelado em Ciência da Computação) - Centro de
Ciências Exatas e Naturais, Universidade Regional de Blumenau, Blumenau.

\noindent\textbf{[verbete de enciclopédia em meio eletrônico:]}

EDITORES gráficos. In: WIKIPEDIA, a enciclopédia livre. [S.l.]:
Wikimedia Foundation, 2006. Disponível em:
\url{http://pt.wikipedia.org/wiki/Editores_graficos}. Acesso em: 13 maio
2006.

\noindent\textbf{[artigo em evento:]}

FRALEIGH, Arnold. The Algerian of independence. In: ANNUAL MEETING OF
THE AMERICAN SOCIETY OF INTERNATIONAL LAW, 61, 1967, Washington.
\textbf{Proceedings\ldots{}} Washington: Society of International Law,
1967. p. 6-12.

\noindent\textbf{[norma técnica:]}

IBGE. \textbf{Normas para apresentação tabular}. 3. ed. Rio de Janeiro,
1993. 61 p. Disponível em:
\url{http://biblioteca.ibge.gov.br/visualizacao/monografias/GEBIS%20-%20RJ/normastabular.pdf}.
Acesso em: 27 ago. 2013.

\noindent\textbf{[artigo em periódico:]}

KNUTH, Donald E. Semantic of context-free languages.
\textbf{Mathematical Systems Theory}, New York, v. 2, n. 2, p. 33-50,
jan./mar. 1968.

\noindent\textbf{[trabalho acadêmico ou monografia (TCC/Estágio, especialização, dissertação, tese):]}

SCHUBERT, Lucas A. \textbf{Aplicativo para controle de ferrovia
utilizando processamento em tempo real e redes de Petri}. 2003. 76 f.
Trabalho de Conclusão de Curso (Bacharelado em Ciência da Computação) -
Centro de Ciências Exatas e Naturais, Universidade Regional de Blumenau,
Blumenau.

\noindent\textbf{[página da internet com autor]}

SCHULER, João P. S. \textbf{Tutorial de Delphi}. Porto Alegre,
[2002]. Disponível em: \url{http://www.schulers.com/jpss/pascal/dtut/}.
Acesso em: 27 ago. 2013.

\noindent\textbf{[página da internet sem autor]}

SCHRATCH. \textbf{Program, imagine, share}. [S.l.], [2013?].
Disponível em: \url{https://scratch.mit.edu/}. Acesso em: 27 maio 2013.

\noindent\textbf{[relatório de pesquisa:]}

VARGAS, Douglas N. \textbf{Editor dirigido por sintaxe}. 1992. Relatório
de pesquisa n. 240 arquivado na Pró-Reitoria de Pesquisa, Universidade
Regional de Blumenau, Blumenau.
